\section{Planning}

\beameronly{
\begin{frame}{Task: Software Demo Nov 25th}
    \mytext{
    
    \begin{itemize}
        \item Decide on a Software, state on forum
        \item Make yourself an account
        \item Create a (dummy) project
        \item Enter tasks, Work packages, responsibilities etc.
        \item Create Gantt diagrams, other views
        \item Move a milestone
        \item Demo this to the class
    \end{itemize}
    }
    \mycomment{
    Demos must not use slides, and should be around 15 minutes.
    }
\end{frame}
}

\begin{frame}{No Planning: Bias for Action}
    \mycomment{
    We might think planning is prerequisite for action, but it isn't. Planning is a science, that not everyone is trained at. For many roles, planning is not even required!
    }
    \mytext{
    \begin{itemize}
        \item Tendency to act without plan
        \item Makes sense when there are no risks
        \item Example: business-as-usual
        \item Planning is often under-financed or not financed
        \item Optimism bias: need to plan is under-estimated
    \end{itemize}
    }
\end{frame}

\begin{frame}{Planning}
    \mytext{
    Planning, a genuine human capacity, is to
    \begin{itemize}
        \item define a desired state (the goal),
        \item lay out the steps necessary to achieve it,
        \item lay out the steps necessary to achieve these, etc.,
        \item bring them into an efficient order.
    \end{itemize}
    Plans are based on available information; but the information we have grows with planning, and again with execution. So \textbf{plans change.}
    }
    \mycomment{
    Animals do not plan, or when they plan, only a few steps ahead. Humans have the capacity to make plans: to project a goal in the future, and to take multiple steps to achieve them. 
    We can build, but we can build the tools we need to build, and so forth.
    }
\end{frame}


\begin{frame}{Agenda, to-do-lists}
    \mycomment{
    A popular way to plan things is to write everything that needs to be done down at once---on a \enquote{to-do-list}. While this is a good starting point, such lists have severe shortcomings.
    }
    
    \mytext{
    \begin{itemize}
        \item \emph{agenda} n.pl.~lat.~`(things) to be done'
        \item problem: top priority may delay everything else
        \item things stay on list until they become urgent
        \item possible multi-tasking?
        \item division of labour?
        \item better: define a date (point in time)! $\rightarrow$ \textit{schedule!}
        \item have multiple agenda for different dates
        \item colour scheme (red-yellow-green)
        \item \textbf{software}
    \end{itemize}
    }
\end{frame}


\begin{frame}{\enquote{Eisenhower} Matrix}
    \showmytikz{../shared/tikz/tikz-eisenhower}
    \mycomment{
        The so-called \enquote{Eisenhower} matrix shows a way you could classify tasks from your to-do-list. Such things that are absolutly urgent should be done immediately without hesitation. Every thing else should either be scheduled, delegate, or simply deleted.
    }
\end{frame}


\begin{frame}{Project Management Software}
    \mycomment{
    These considerations should suffice to understand why planning, and project management in general, should not be left to a sheet of paper, or a paper calendar. Today, we will use Project Management Software. 
    }
    \mytext{
    \begin{itemize}
        \item Why?
        \item What to look for?
        \item Criteria?
    \end{itemize}
    }
\end{frame}


\begin{frame}{Foresight, Planning}
    \mytext{
    \begin{itemize}
        \item Efficiency comes with planning. 
        \item \textbf{Planning dominates the management of complex  projects.}
        \item Complexity may not scale linearly.
        \item Planning is a communicative activity,
        \item Planning has to do with foresight
        \item \textit{What must be done?}
        \item \textit{And in which order?}
        \item \textit{By whom, with what?}
        \item \textit{Depending on what?}
        \item What are the sub-projects or tasks?
        \item \textbf{Plans change}
    \end{itemize}
    }
\end{frame}

% \begin{frame}{Objective and Owner of a Project}
%     \begin{itemize}
%         \item \textbf{Who }wants what?
%         \item Project owner
%         \item Project sponsor
%         \item Project governance
%         \item Common vision?
%         \item Well defined?
%     \end{itemize}
% \end{frame}

% \begin{frame}{Project Flow Chart (NEW)}
%     \showmytikz{../shared/tikz/tikz-projectflowchartfour}
    
% \end{frame}

\begin{frame}{Project Flow Chart}
    \mycomment{
    Most projects can at best be understood in form of a flow chart to demonstrate its temporal order.
    }
    \mytext{
    \begin{itemize}
        \item X initialised the project
        \item to achieve Y;
        \item Business environment = project began
        \begin{itemize}
            \item planning
            \item execution
            \item closure
        \end{itemize}
        \item X was achieved / running
    \end{itemize}
    }
\end{frame}

\begin{frame}{Project Flow Chart}    
    \mycomment{
    In the simplest form, such a flow chart will correspond to the typical phases a project has.
    }
    \showmytikz{../shared/tikz/tikz-PM-project-flow-linear}
    \end{frame}

\begin{frame}{Project Flow Chart}
    \mycomment{
    The parts from the flow diagram above can be re-written: the supervision and governance exists throughout the entire project; then, there is the temporal axis of planning, execution and final product. The end result will have to be justified before the project sponsor, who is at the origin of the project.
    }
    \showmytikz{../shared/tikz/tikz-PM-project-flow-branching}
    \end{frame}

\begin{frame}{Project Flow Chart}    
    \mycomment{
    The temporal axis can be further devided: there will be multiple  rounds and feedback loops in the planning phase; the execution will consist of many sub-projects; and the final product will have multiple dimensions or qualities.
    }
    \showmytikz{../shared/tikz/tikz-PM-project-flow-branching-labeled}
    
    
\end{frame}

\begin{frame}{What needs to be done?}
    \mycomment{
    When we want to start planning, we must identify the requirements and the required tasks first; then, identify the resources each requires, and their temporal order.

    But most importantly, we must understand that every planning is preliminary. Even the best of plan will have to change for unforeseen reasons. Planning must take incertainty into account.
    }
    \mytext{
    \begin{itemize}
        \item \textbf{Planning starts with listing tasks}
        \item[$\Rightarrow$] \textbf{and defining their deliverables}
        \item \textbf{Planning is scheduling tasks}
        \item[$\Rightarrow$] \textbf{and allocating resources}
        \item \textbf{Plans change!}
    \end{itemize}
    }
\end{frame}

\begin{frame}{Sub-Projects are called 'Work Packages'}
    \mycomment{
    As discussed, projects consists of sub-projects. As these are assigned to responsible parties, they are called \enquote{work packages}. 
    }
    
    \mytext{
    
    Sub-projects are often distributed to different responsible parties, hence the name: 
    
    
    \textbf{Work package} (WP)
    
    % \medskip
    
    % German: `Arbeitspaket (AP)' - `Teilarbeitspaket'\\
    % English: Subproject, Task, Work Package
    }
\end{frame}


\begin{frame}{Break Work down to Tasks}
    \mycomment{
    At first, identify task carefully; group them by work package.
    }
  \mirror{\showmytree[1]{../shared/tikz/forest-work-breakdown.tex}}
\end{frame}


\begin{frame}{Break Work down to Tasks}
    \mycomment{
    The \enquote{normal} way to present this is the other way round.
    }
  \showmytree[1]{../shared/tikz/forest-work-breakdown.tex}
\end{frame}



\begin{frame}{Break Work down to Tasks}
    \mycomment{Then, assign these work packages to someone responsible. In larger projects, these would be consortium partners or contractors. In smaller projects, team leaders. }
  \showmytree[2]{../shared/tikz/forest-work-breakdown.tex}
\end{frame}

\begin{frame}{Break Work down to Tasks}
    \mycomment{
    In the next step, each task will be assigned a responsible person. Notice that here, we have identical persons working in multiple work packages. If this is possible or not of course depends on the actual governance structure. It is, however, advisable to cross-assign responsibilities, so that individual task remain connected with the larger picture. 
    }
  \showmytree[3]{../shared/tikz/forest-work-breakdown.tex}
\end{frame}

\begin{frame}{Break Work down to Tasks}
    \mycomment{
    Finally, each of these tasks must be assigned temporal and financial and other resources. This can be estimated at the beginning, but once there is a responsible person, this person must be given a say on the assumed duration, the requirements, and the resources needed.
    }
  \showmytree[4]{../shared/tikz/forest-work-breakdown.tex}
\end{frame}


% \input{TE_Hatfactory}

\begin{frame}{Flow Charts, Lists, Tree Chart, Gantt Diagramm --}
    
    \mytext{
        The Work Breakdown diagram in form of a tree is a helpful starting point. But now we realise that there is a temporal order of these tasks: some can be worked on in parallel; some depend on one another so that one can only start once the other is completed; and we would like to sum up total resource requirements. So we need other forms to visualise the same project data!
    }
\end{frame}

\begin{frame}{Flow Charts, Lists, Tree Chart, Gantt Diagramm --}
    
    \mytext{
    You need an overview of what is to be done!
    
    The overview can take the form
    
    \medskip

    \begin{itemize}
        \item \colorbox{green!40}{flow chart},a.k.a. \textit{network diagram}
        \item \colorbox{lightgray}{List}, a.k.a. \textit{register}
        \item \colorbox{yellow!40}{WBS}, a.k.a. \textit{tree diagram}
        \item \colorbox{red!40}{Gantt Diagram}
    \end{itemize}
    }
\end{frame}

\begin{frame}{Flow Charts, Lists, Tree Chart, Gantt Diagramm---}
    \mytext{
    You always need all three forms!
    
    \showmytikz{../shared/tikz/tikz-wbs-tree-gantt-list}
    }
    \mycomment{
        Project management software should give us the possibility to look at the entire project (or parts of it) in four principal, different forms. 
       
    }
\end{frame}


\begin{frame}{Flow Chart}
    \mytext{
        A \colorbox{green!40}{\textbf{flow chart}} makes the temporal order and dependencies clear, but does not provide much more information than that.
    \showmytikz{../shared/tikz/tikz-chain-kitchen}
    }
    \mycomment{

    Consider the simple project to renovate a kitchen. We see at once three phases, or work packages.

    This diagram shows very clearly what has to be done in which order. But it does not show how long each part will take, how much it will cost, who is in charge, etc.   
    
    }
\end{frame}

\begin{frame}{List or Register}
    \mytext{
    The most comprehensive form is a \colorbox{lightgray}{\textbf{list / register}}. However, as columns become very large, this can become difficult to read. Also, dependencies and temporal order are not visible. 
    
    \medskip
    
    \mytable{\scriptsize
        \begin{tabular}{rllrrllr}
        \toprule
        & Process & Result & hrs. & Dep. & Assigned & Resources & Budget\\
        \midrule
        1. & Manage & (the project) & cont. & management & Alice & Telephone & 100 \\
        2. & Empty room & room empty & 1 & movers & Bob & space & 100\\
        3. & Paint room & paint dry & 2 & painters & Charlie & paint etc.&  100 \\
        4. & Change floor & new floor & 2 & floorlayers & Bob & floorboards & 300 \\
        5. & Clear out & rubbish disposed & 2 & movers & Bob & van & 100 \\
        6. & Replace items & new appliancies & 2 & management & Alice & suppliers  & 4000 \\
        7. & Refurbish & livable room & 2 & Users? & Alice & movers & 500 \\
  
        \bottomrule
        \end{tabular}
        }
    }
\end{frame}

\begin{frame}{Work Breakdown Structure}
    \showmytikz{../shared/tikz/forest-houshold-wbs}
    \mycomment{
       The \colorbox{yellow!40}{\textbf{WBS Work Breakdown Structure}} (in form of a tree) is very informing, but notoriously difficult to produce. But you see that the contents of the WBS and the contents of the list are indentical. A good software solution will allow us to switch from one view to the other. 
    }
\end{frame}

\begin{frame}{Gantt Diagram}
    \mytext{
    \colorbox{red!40}{\textbf{Gantt diagram}} is the most common visualisations of projects, which show very clearly duration and dependencies. 

    \medskip
    
    \showmytikz{../shared/tikz/tikz-PM-gantt-kitchen-renovation}
    }
    \mycomment{
        They do not, however, provide much detail, e.g. about the resources requred or persons in charge. 
        
        They are also notoriously difficult to draw and to keep up to date---unless you use appropriate software, that is. 
        
    }
\end{frame}


% \input{TE_Hatfactory}



% \begin{frame}{Example for Work Breakdown Structure}
%   \begin{forest}
%     for tree={
%         font=\footnotesize,
%         grow = east,
%         anchor=west,
%         forked edge,
%         inner sep=2pt
%             }
%     [\textbf{Have a new ship}
%     ]
%     \end{forest}
    
%     \framebreak
    
%     \begin{forest}
%     for tree={
%         font=\footnotesize,
%         grow = east,
%         anchor=west,
%         forked edge,
%         inner sep=2pt
%             }
%     [\textbf{Have a new ship}
%         [SP3 use ship]
%         [SP2 order ship]
%         [SP1 plan ship]
%     ]
%     \end{forest}    
    
%     \framebreak
    
%     \begin{forest}
%     for tree={
%         font=\footnotesize,
%         grow = east,
%         anchor = west,
%         forked edge,
%         inner sep=2pt
%             }
%     [\textbf{Have a new ship}
%         [SP3 use ship [find port][find jobs for ship][find crew]]
%         [SP2 order ship [supervise][contract][find contractor]]
%         [SP1 plan ship [technical planning][business planning]]
%     ]
%     \end{forest}        

% \end{frame}


\begin{frame}{Data Structure of Tasks}
    \mytext{
    We have seen that all four forms: \colorbox{green!40}{flow chart}, \colorbox{yellow!40}{tree chart (WBS)}, \colorbox{red!40}{Gantt chart} and \colorbox{lightgray}{list form} are different visualisation of the same core data. 

    
    \mytable{
    \begin{tabular}{l l l l l l}
    \toprule
         Task & Result &  Resp. & Duration & Budget & Depend. \\
         \midrule
         1 & 1.1 & Doreen & 9d & 5K & -\\
         2 & 2.1 & Alice & 4d & 2K & 1.1 \\
           & ... & ... & ... & ... & ... \\
         \bottomrule
    \end{tabular}
    }
    
        Every task is a data entry. The underlying structure should be a relational database, in that there are clearly defined dependencies between the different fields. 
        
    }
    \mycomment{
        For example, there must be a data base with personal data (of Doreen, Alice and so forth), and each task is assigned to exactly one such person. 


    }

    \end{frame}

\begin{frame}{Data Structure of Tasks}
    \mycomment{
            This shows why project management software are database applications. Running a word document or an Excel sheet will not do the job.    
    }
    \mytext{
    
    Every task is an entry in a database with...
    
     \begin{itemize}
         \item Area, or 'Work package'
         \item Task number
         \item Task name
         \item Duration
         \item Result, or Milestone
         \item Dependencies (backward, forward)
         \item Responsibility
         \item Budget
         \item Other ressources
     \end{itemize}
    }
\end{frame}

\begin{frame}{Data Structure of Tasks}
    
    \mytext{
    Every task is an entry in a database with...
    
     \begin{itemize}
         \item Area, or 'Work package'
         \item Task number
         \item Task name
         \item Duration
         \item Result, or Milestone
         \item Dependencies (backward, forward)
         \begin{itemize}
             \item Responsibility
             \item Budget
             \item Other ressources
         \end{itemize}
     \end{itemize}
         
    
    }
\end{frame}


\begin{frame}{From WBS to Gantt-Diagram}
    \mycomment{
    Plans cannot be perfect. There are some rules to follow: one, \textbf{each task results in a milestone,} and this milestone is the requirement for the dependent task.

    Perfect planning takes uncertainty into account, and does so by building in buffers. For example, builders need a break; people need a holiday; machines have a maximum time they can run; work can only be done in daytime; deliverances may come too late; additional paperwork might be needed. 
    }
    
    \mytext{
    
    \begin{itemize}
        \item Every task should end in a milestone
        \item Define dependencies!
        \item Estimate time (carefully)
        \item Build in buffers
        \item Calculate resources
        \begin{itemize}
            \item no idle time
            \item no overlaps
        \end{itemize}
        \item Use visualisation / software
        % \item[] $\rightarrow$ \textbf{reasearch!}
    \end{itemize}
    
    }
\end{frame}



\begin{frame}{Data Structure of Tasks}
    
    \mytext{
    
     IMPORTANT: 
     
     \begin{itemize}
         \item Every \textbf{task} or \textbf{work} is a \textbf{process!}
         \item[] \textit{in contrast}
        \item every \textbf{result} is a \textbf{milestone!} 
        \begin{itemize}
            \item a milestone can only be what is a result!
            \item processes are not results and not milestones!
            \item milestones are not processes!
        \end{itemize}

     \end{itemize}
    }
\end{frame}

\begin{frame}{Ideal PM Software Solution}
    
    \mytext{
    
    
    We seek a software that
    
    \begin{enumerate}
        \item Allows us to mind-map \colorbox{yellow!40}{WBS}
        \item turns this into a database (\colorbox{lightgray}{list})
        \item shows this as \colorbox{red!40}{Gantt diagram}
        \item allows us to \textbf{move things} 
        \item allows us to \textbf{record changes}
        \item allows us to \textbf{monitor progress}
        \item allows us to \textbf{monitor costs}
        \item allows \textbf{team use} it
        \item allows \textbf{export into presentable forms}
    \end{enumerate}
    
    }
\end{frame}



\begingroup
    \ifdefined\setbeamercolor\setbeamercolor{background canvas}{bg=violet!20}\fi
    \begin{frame}{TASK}
    
    \mytext{
    \begin{enumerate}
        \item Choose a (small) example project
        \item find a Project Management software
        \item create a free account
        \item create a project
        \item brainstorm the different tasks
        \item enter these tasks
        \item create a Gantt-Diagram
        \item create dependencies between tasks
        \item move the end date of the first task
        \item \textbf{this you demonstrate to the class}
    \end{enumerate}
    }
    \end{frame}

\endgroup
    
